\documentclass{easychair}

\usepackage{amsmath,amssymb,amsthm}

\newtheorem{definition}{Definition}
\newtheorem{theorem}{Theorem}
\newtheorem{proposition}{Proposition}
\newtheorem{lemma}{Lemma}

\title{Fuzzifying Derivatives of Regular Expressions\\
\normalsize{(Extended Abstract)}}

\author{
Mariam Katamashvili\ \ \ \ \  David Maisuradze
}

\institute{Kutaisi International University, Kutaisi, Georgia}

%\authorrunning{Mariam Katamashvili and David Maisuradze}
%\titlerunning{Derivative of Regular Expressions with Similarity Relation}

\begin{document}

\maketitle
There is growing interest in the Computer Science and AI communities in moving from {\em exact},  {\em crisp}, {\em qualitative} settings to {\em metric}, {\em probabilistic}, {\em quantitative} and hence  ultimately {\em fuzzy} settings. In this view, we explore how to extend the standard theory of finite automata~\cite{C71}, based on exact equality of symbols, to a setting where symbols are fuzzy~\cite{DHMM}. Namely, we assume that the alphabet is enriched with a  {\em similarity
relation} and matching is performed up to {\em cut} values. In this extended abstract we generalize the theory of  
Brzozowski~\cite{B62} derivatives, which  gives a direct construction of automata from regular
expressions: after reading a word, the current state is represented by the
corresponding derivative of the expression. All our results are formalized in Rocq~\cite{Rocq}.

Let $\Sigma$ be a finite alphabet and let $
\mathcal R:\Sigma\times\Sigma\to[0,1]$
be a {\em fuzzy similarity relation}, where $[0,1]$ is endowed with a left-continuous T-norm, which makes it into a residuated lattice~\cite{met05}. Throughout this extended abstract, we use the minimum T-norm for simplicity. A similarity relation satisfies reflexivity, symmetry and the reverse triangle inequality.
%For $0<\mu\le1$, its $\mu$-cut is
%\[\mathcal R_\mu=\{(a,b)\in\Sigma\times\Sigma\mid \mathcal R(a,b)\ge\mu\}.\]

Similarity is extended from symbols to words, languages, and regular expressions. 
\begin{definition}\hfill\\
$\bullet$ {\bf Similarity on words.}
For $w_1,w_2\in\Sigma^\ast$, define $
\mathcal R^w:\Sigma^\ast\times\Sigma^\ast\to[0,1]$
recursively by\\ $
\mathcal R^w(w_1,w_2)=
\begin{cases}
1 & \text{if } w_1=\varepsilon \text{ and } w_2=\varepsilon,\\[2pt]
0 & \text{if } (w_1=\varepsilon \text{ and } w_2\neq\varepsilon)
      \text{ or } (w_2=\varepsilon \text{ and } w_1\neq\varepsilon),\\[4pt]
\min\!\bigl(\mathcal R(a_1,a_2),\mathcal R^w(w_1',w_2')\bigr)
& \text{if } w_1=a_1w_1' \text{ and } w_2=a_2w_2'.
\end{cases}$\smallskip\\
$\bullet$ {\bf Similarity on languages}
For languages $L_1,L_2\subseteq\Sigma^\ast$, define
\[
\mathcal R^L(L_1,L_2)
=
\min\!\Bigl\{
\inf_{w_1\in L_1}\sup_{w_2\in L_2}\mathcal R^w(w_1,w_2),
\inf_{w_2\in L_2}\sup_{w_1\in L_1}\mathcal R^w(w_1,w_2)
\Bigr\}.
\]\end{definition} We do not spell out the definition of similarity $\mathcal R^r$ on regular
expressions, for lack of space. This definition is intentionally syntax-directed and therefore
syntax-sensitive. Constructors such as intersection and complement do not preserve the
desired syntax-to-semantics inequalities in general, so they are not amenable to a well-behaved fuzzification. Hence, the theorems are restricted to the fragment generated by $\emptyset,\varepsilon$,
letters, union, concatenation, and Kleene star. 

\begin{proposition} The relations $\mathcal R^w$, $\mathcal R^L$, and $\mathcal R^r$ are similarity relations. 
\end{proposition}

One of the main results of this work is that of establishing a connection between our fuzzy analogue of the theory of finite automata and the crisp version of the theory for a corresponding quotient language. Namely, we introduce the following definitions.
\begin{definition}For a fixed cut value $\mu$, define the  following equivalence relations: \\
$\bullet$ $\boldsymbol{\sim_\mu^\Sigma}$ on the alphabet $\Sigma$ by

$a\sim_\mu^\Sigma b
\quad\Longleftrightarrow\quad
\mathcal R(a,b)\ge\mu;
$\\
$\bullet$ $\boldsymbol{\sim_\mu^{\Sigma^\ast}}$ on words in $\Sigma^\ast$ by 

$w_1\sim_\mu^{\Sigma^\ast} w_2
\quad\Longleftrightarrow\quad
\mathcal R^w(w_1,w_2)\ge\mu;
$\\
$\bullet$ $\boldsymbol{\sim_\mu^{\mathcal P(\Sigma^\ast)}}$ on languages over $\Sigma^\ast$ by 

$L_1\sim_\mu^{\mathcal P(\Sigma^\ast)} L_2
\quad\Longleftrightarrow\quad
\mathcal R^L(L_1,L_2)\ge\mu.
$
\end{definition}
We can now introduce the appropriate quotient notions.
\begin{definition}Let
$\Sigma^\mu=\Sigma/{\sim_\mu^\Sigma}$
be the {\em quotient alphabet} with quotient map
$q_\mu:\Sigma\to\Sigma^\mu$ denoted by $q_\mu(a)=[a]_\mu$, and define:\\
$\bullet$ the extension of $q_\mu$ to words by putting $
q_\mu^\ast(a_1\cdots a_n)
=
[a_1]_\mu\cdots[a_n]_\mu
$ ; \\
$\bullet$ the  quotient image  of $L\subseteq\Sigma^\ast$ by putting
$ L^\mu=%q_\mu^\ast(L) =
\{q_\mu^\ast(w)\mid w\in L\}
\subseteq(\Sigma^\mu)^\ast.
$ 
\end{definition}
Finally, we give the main definitions.\begin{definition}[Fuzzy language derivative]
For $L\subseteq\Sigma^\ast$ and $a\in\Sigma$, define\\
\ \ \ \ $
\partial_a^\mu(L)
=
\{w\in\Sigma^\ast\mid
\exists b\in\Sigma,\exists v\in\Sigma^\ast:
\mathcal R(a,b)\ge\mu,\ 
\mathcal R^w(w,v)\ge\mu,\ 
bv\in L
\}.
$
\end{definition}
\begin{definition}[Fuzzy expression derivative]
For a regular expression $r$, the syntactic fuzzy derivative is
$
\mathcal D_a^\mu(r)
=
\operatorname{Sat}_\mu
\left(
\sum_{b\in S_a^\mu}\mathcal D_b(r)
\right)
$
where $\mathcal D_b(r)$ is the ordinary Brzozowski derivative, and for  $L\subseteq\Sigma^\ast$, the $\mu$-saturation is given by $
\operatorname{Sat}_\mu(L)
=
\{w\in\Sigma^\ast\mid
\exists v\in L \text{ such that } w\sim_\mu^{\Sigma^\ast}v\}.
$
\end{definition}

We are now in the position of giving the main results. 

\begin{proposition} For all (restricted) regular expressions $r_1,r_2$,\\
$\bullet$ 
$
\mathcal{R}^r(r_1,r_2)
\le
\mathcal{R}^L\bigl(L(r_1),L(r_2)\bigr).
$\\
$\bullet$ $\mathcal{R}^L\bigl(L(r),L(s)\bigr)
=
\sup\Bigl\{
\min\bigl(
\mathcal{R}^r(r',t),
\mathcal{R}^r(s',u)
\bigr)
\;\Bigm|\;% r',s',t,u \text{ are classical regular expressions}, \;
L(r')=L(r),
L(s')=L(s),
L(t)=L(u)
\Bigr\}.$
\end{proposition}


\begin{theorem}
$
\partial_a^\mu(L)
=
\{w\in\Sigma^\ast\mid
q_\mu^\ast(w)\in
\partial_{[a]_\mu}(L^\mu)\}.
$
\end{theorem}
Thus, the fuzzy derivative over $\Sigma$ can be understood as the ordinary
derivative over the quotient alphabet $\Sigma^\mu$, then transferred back to
words over $\Sigma$.
\begin{theorem}[Regularity preservation]
If $L\subseteq\Sigma^\ast$ is regular, then
$\partial_u^\mu(L)$
is regular for every $u\in\Sigma^\ast$ and every $0<\mu\le1$.
\end{theorem}

\begin{theorem}[Semantic--syntactic correspondence]
For every regular expression $r$, word $u\in\Sigma^\ast$, and cut value
$0<\mu\le1$, $
L(\mathcal D_u^\mu(r))
=
\partial_u^\mu(L(r)).
$
\end{theorem}

\begin{theorem}[Finiteness]\label{finiteness}
If $\Sigma$ is finite and $L\subseteq\Sigma^\ast$ is regular, then
$
\{\partial_u^\mu(L)\mid u\in\Sigma^\ast\}
$
is finite for every fixed $0<\mu\le1$.
\end{theorem}

In conclusion, fuzzy automata can be viewed as quotients of ordinary derivative
automata over the quotient alphabet $\Sigma^\mu$. 
\begin{theorem}The language recognized by
the resulting quotient automaton is the concrete $\mu$-saturation{\color{red}??}, namely
$
\mathcal L_\Sigma^\mu(r)
=
(q_\mu^\ast)^{-1}(q_\mu^\ast(L(r))).
$
\end{theorem}

All our results are formalized in Rocq~\cite{Rocq,RocqRegexpBrzozowski}, see~\cite{OurRocqRepo}. We build on the  Coquand--Siles~\cite{CoquandRegexEquality} formalization of
regular-expression equivalence via Brzozowski derivatives. Theorem~\ref{finiteness} is crucial in that all notions of finiteness are not equivalent in a constructive setting. We are currently working on proving a Kuratowski-finiteness result by working with normalized representatives of regular
expressions, so that repeated derivatives can be checked to generate only finitely many distinct
states using the equivalence procedure in~\cite{CoquandRegexEquality}. This should be a slight improvement in that the result in~\cite{CoquandRegexEquality} is based on a weaker notion of finiteness.

%{\small \paragraph{Acknowledgments.}
%We would like to thank Furio Honsell and Besik Dundua, for their guidance and support.}

\bibliographystyle{plain}

\begin{thebibliography}{10}

\bibitem{B62}
Janusz A. Brzozowski.
\newblock Canonical regular expressions and minimal state graphs for definite events.
\newblock \emph{Mathematical Theory of Automata}, 12:529--561, 1962.

\bibitem{C71}
John H.Conway
\newblock {\em Regular Algebra and Finite Machines}.
\newblock {Chapman and Hall (London)}, 1971.
 


\bibitem{CoquandRegexEquality}
Thierry Coquand and Vincent Siles.
\newblock A Decision Procedure for Regular Expression Equivalence in Type Theory.
\newblock In \emph{Proceedings of the 1st International Conference on Certified Programs and Proofs (CPP)}, 2011.

\bibitem{DHMM}
Besik Dundua, Furio Honsell, Mariam Katamashvili, and David Maisuradze.
\newblock{\em Fuzzy Automata} (to appear).


\bibitem{OurRocqRepo}
Mariam Katamashvili and David Maisuradze.
\newblock Rocq formalization for regular-expression derivatives with similarity.
\newblock \texttt{https://github.com/DUT1K1/regular-expression-derivative}.

\bibitem{met05}
{George Metcalfe},
\newblock{Fundamentals of Fuzzy Logic},
\newblock \texttt{https://api.semanticscholar.org/CorpusID 17549644}.


\bibitem{OwensReppyTuron}
Scott Owens, John Reppy, and Aaron Turon.
\newblock Regular-expression derivatives re-examined.
\newblock \emph{Journal of Functional Programming}, 19(2):173--190, 2009.

\bibitem{Rocq}
The Rocq Prover Development Team.
\newblock {The Rocq Prover (Version x.x)}, 202x.\\
\newblock \texttt{https://rocq-prover.org}.

\bibitem{RocqRegexpBrzozowski}
Rocq community.
\newblock \emph{regexp-Brzozowski}.
\newblock \texttt{https://github.com/rocq-community/regexp-Brzozowski}.

\end{thebibliography}

\end{document}